\section{Related Work}
\label{sec:related}

\subsection{Spurious Correlations and Worst-Group Accuracy}

Neural networks trained with ERM exploit spurious correlations --- dataset artifacts where simple co-occurring features predict class labels during training but fail to generalize across subpopulations.
On Waterbirds \citep{Wah2011Caltech,Sagawa2020Distributionally} and CelebA \citep{Liu2015Deep}, spurious correlations between background/hair color and class labels cause ERM models to achieve high average accuracy (97\% and 95\% respectively) while failing catastrophically on minority groups (worst-group accuracy 72\% and 47\%) \citep{Sagawa2020Distributionally}.
The theoretical basis is established: \citet{Shah2020Pitfalls} formally prove simplicity bias; \citet{Arpit2017Closer} show DNNs memorize easy-to-fit patterns before hard ones, producing temporal structure exploitable for minority identification.

\subsection{Label-Free Debiasing Methods}

Methods requiring group annotations (GroupDRO~\citep{Sagawa2020Distributionally}) achieve strong worst-group performance but assume the label structure one wishes to avoid.
In the annotation-free regime: JTT~\citep{Liu2021Just} identifies minority as misclassified samples in an initial ERM run (achieving 86.7\% worst-group on Waterbirds under their experimental setup); LfF~\citep{Nam2020Learning} uses relative per-sample loss online.
Both succeed because per-sample loss is an effective proxy --- minority samples incur high loss where the spurious feature conflicts with the true label.
However, this magnitude-based proxy conflates hard-majority samples with spurious-minority samples \citep{Liu2021Just,Kirichenko2022Last}.
DFR~\citep{Kirichenko2022Last} achieves near-supervised performance by retraining a linear head on a balanced held-out set, but requires a labelled validation split.

Our work extends the annotation-free regime to gradient \emph{direction} as a proxy signal, testing whether directional information distinguishes spurious-minority from hard-majority samples in a way loss magnitude cannot.
We note that our contribution is signal existence (ROC-AUC measurement), not downstream task performance, and we do not reproduce JTT or GroupDRO baselines in our experimental setup.

\subsection{Gradient-Based Training Dynamics Analysis}

Per-sample gradient analysis has a growing role in understanding training.
\citet{Katharopoulos2018Not} use gradient norms for importance sampling.
\citet{Toneva2019Empirical} track forgetting events.
\citet{Swayamdipta2020Dataset} characterize data cartography using per-epoch loss and correctness signals.
Most relevant: \citet{BiasTrail2025Bias} demonstrates gradient probes computed at layers can identify biased samples without labels using directional gradient signals with globally stable references --- consistent with our finding that per-batch reference is the failure point.
Gradient alignment has also been studied in continual learning (GEM~\citep{LopezPaz2017Gradient}) and multi-task learning~\citep{Yu2020Gradient}, where global or cross-task reference directions are used --- not per-mini-batch statistics.

\textbf{Our contribution:} No prior study has directly measured within-batch gradient alignment ROC-AUC as a spurious-minority detector on Waterbirds or CelebA across training epochs.
We fill this gap, characterize the failure mode, and identify the reference direction correction required.
